
\svnid{$Id: $}
\svnid{$URL: $}

\section{FOR REVIEW:junction advection for a flood wave (flow confluence and acceleration at junction; equidistant grid; positive flow upstream) applying conservation of momentum (iadvec1d=1)}
\newrefsegment

%---------------------------------------------------------------------------------
\section*{Quality Assurance}
\begin{tabular}{@{}p{27mm}@{}p{0mm}p{\textwidth-52mm-48pt}p{15mm}p{10mm}}
\textbf{Date} & \textbf{:} &   May 16 2019\\ 
\textbf{Author} & \textbf{:} & Christianne Luger  & \textbf{Initials:} & CL \\
\textbf{Review} & \textbf{:} &   & \textbf{Initials:} & \ldots 
\end{tabular}

%---------------------------------------------------------------------------------
\section*{Version information}
{{
\begin{tabular}{@{}p{27mm}@{}p{0mm}p{\textwidth-27mm-24pt}}
\textbf{Executable}   & \textbf{:} & Deltares, D-Flow FM Version 1.2.54.64224M, Jun 28 2019, 15:14:31  \\
\textbf{Location input}     & \textbf{:} & \url{https://repos.deltares.nl/repos/DSCTestbench/trunk/cases/e106_dflow1d/f04_numerical-aspects/c14_junction_advection08_iadvec1d_1} \\
\textbf{Revision input} & \textbf{:} & SVN revision number of input \\
\textbf{Location output}     & \textbf{:} & \url{https://repos.deltares.nl/repos/DSCTestbench/trunk/cases/e02_dflowfm/f104_1D_numerical-aspects/c07_junction-advection-flood-wave-confluence-equidistant-pos-flow-upstream} \\
\textbf{Revision output} & \textbf{:} & SVN revision number of reference output
\end{tabular}
}}

%---------------------------------------------------------------------------------
\section*{Purpose}
\textbf{FOR REVIEW} 
The purpose of test \textbf{e02}(dflowfm)\_\textbf{f104}(numerical-aspects)\_\textbf{c07}(junction-advection-flood-wave-confluence-equidistant-pos-flow-upstream) is to verify:

\begin{enumerate}
	\item the implementation of junction advection (see \autoref{10407:description of junction advection}) for a flood wave (flow confluence and acceleration at junction; equidistant grid; positive for upstream) while applying conservation of momentum (\textbf{\textcolor[rgb]{1,0,0}{iadvec1d=1}}).
	\item that the time-varying discharge and constant water level boundaries work. 

\end{enumerate}

%---------------------------------------------------------------------------------
\section*{Linked claims}
\begin{enumerate}
	\item Junction advection is correctly implemented for a flood wave (flow confluence and acceleration at junction; equidistant grid; positive for upstream) while applying conservation of momentum (iadvec1d=1). 
	\item Time-varying (upstream) discharge boundaries are correctly applied on the systems.
	\item Constant (downstream) water level boundaries are correctly applied on the systems.
\end{enumerate}


%---------------------------------------------------------------------------------
\section*{Approach}

\begin{enumerate}

	\item Water levels and velocities near junctions are compared with the results of identical models having no junctions (or connection nodes). More precisely, results of sub-models \textbf{\textcolor[rgb]{1,0,0}{T3}} are validated against the results of sub-models \textbf{\textcolor[rgb]{1,0,0}{T1}} to \textbf{\textcolor[rgb]{1,0,0}{T2}} (see \autoref{fig:10407:Test model lay-out}).
	\item The results are compared to an identical test case. The SOBEK3 model and output used for this comparison can be found on \url{https://repos.deltares.nl/repos/DSCTestbench/trunk/cases/e02_dflowfm/f104_1D_numerical-aspects/c07_junction-advection-flood-wave-confluence-equidistant-pos-flow-upstream/analysis/SOBEK_reference_case_c14_equidistant}. This model was rand through DIMR as the computation core of SOBEK3 no longer supports a rectangular cross section requried for this case. 
\end{enumerate}

\begin{figure}[H]
    \centering
    \includegraphics*[width=0.95\textwidth]{figures/e02_f104_c07_layout.png}
    \caption{Lay-out of the test model, non-equidistant grid spacing at connection nodes}
    \label{fig:10407:Test model lay-out}
\end{figure} 
%---------------------------------------------------------------------------------
\section*{Conclusion} 
The implementation of junction advection for a flood wave (flow confluence and acceleration at junction; equidistant grid; ; positive upstream flow direction) for conservation of momentum (iadvec1d=1) is considered to be most satisfactory. This conclusion is supported by the findings that the water level and velocities in the sub-model with and without a junction (or connection) node are almost identical.

A comparison to an identical SOBEK3 test case show insignificant differences in the water level and velocities in all sub-models. This further supports the conclusion above.

%---------------------------------------------------------------------------------
\section*{Model description} \label{10407:model lay-out}

The test comprises of six separate sub-models \textbf{\textcolor[rgb]{1,0,0}{T1}} to \textbf{\textcolor[rgb]{1,0,0}{T6}} (see \autoref{fig:10407:Test model lay-out}). The upstream boundary condition for all models is a \textbf{time-varying} discharge shown on \autoref{fig:10407:Upstream discharge hydrograph.jpg}. The downstream boundary condition for all models is a \textbf{constant} water level detailed on \autoref{10407:Boundary conditions}. The branch properties, hydraulic roughness and cross-sectional profiles of each sub-model are defined \autoref{10407:Branch direction and h-pnt_coordinates} to \autoref{10407:Open Y-Z rectangular cross-sections}.

\textbf{Please note} that, from the upstream discharge boundary up to and including the downstream water level boundary, per metre of cross-sectional width yields that the hydraulic forcing, the bathymetry and roughness of sub-models \textbf{\textcolor[rgb]{1,0,0}{T3}} to \textbf{\textcolor[rgb]{1,0,0}{T6}} are identical to those of sub-models \textbf{\textcolor[rgb]{1,0,0}{T1}} and \textbf{\textcolor[rgb]{1,0,0}{T2}}.


\textbf{For information only:} \newline
The only difference between test \textbf{e02}(dflowfm)\_\textbf{f104}(numerical-aspects)\_\textbf{c07}(junction-advection-flood-wave-confluence-equidistant-pos-flow-upstream) described in this document and test \textbf{e02}(dflowfm)\_\textbf{f104}(numerical-aspects)\_\textbf{c11}(junction-advection-flood-wave-confluence-non-equidistant-pos-flow-upstream) is the uniformity of grid spacing.

\begin{figure}[H]
    \centering
    \includegraphics*[width=0.48\textwidth]{figures/e106_f04_c14_Upstream_hydrograph.jpg}
    \caption{Upstream discharge per metre of cross-sectional width, that is imposed at the Q-bnd boundaries (see \autoref{fig:10407:Test model lay-out}) of sub-models \textcolor[rgb]{1,0,0}{\textbf{T1}} to \textcolor[rgb]{1,0,0}{\textbf{T6}}.}
    \label{fig:10407:Upstream discharge hydrograph.jpg}
\end{figure}  


\begin{longtable}{|p{25.5mm}|p{17.5mm}|p{17.5mm}|p{17.5mm}|p{17.5mm}|p{17.5mm}|p{17.5mm}|}
\caption{Branches (including branch length, branch direction and coordinates of water level calculation points (h-pnt)) in sub-models \textcolor[rgb]{1,0,0}{\textbf{T1}} to \textcolor[rgb]{1,0,0}{\textbf{T6}}.} 
\label{10407:Branch direction and h-pnt_coordinates} \\ 
\hline
& \begin{center} \textbf{\textcolor[rgb]{1,0,0}{T1}} \end{center}
& \begin{center} \textbf{\textcolor[rgb]{1,0,0}{T2}} \end{center}
& \begin{center} \textbf{\textcolor[rgb]{1,0,0}{T3}} \end{center}
& \begin{center} \textbf{\textcolor[rgb]{1,0,0}{T4}} \end{center}
& \begin{center} \textbf{\textcolor[rgb]{1,0,0}{T5}} \end{center}
& \begin{center} \textbf{\textcolor[rgb]{1,0,0}{T6}} \end{center} \\
  [1ex] \hline \endhead
\hline
\endfoot
%
\STRUT  
\scriptsize{\textbf{B1}: length \unitbrackets{m}}       & \scriptsize{25000.00}     & \scriptsize{25000.00}     & \scriptsize{15000.00}     & \scriptsize{15000.00}     & \scriptsize{15000.00}     & \scriptsize{15000.00}     \\
\scriptsize{\textbf{B1}: direction}                     & \scriptsize{West to East} & \scriptsize{East to West} & \scriptsize{To ConNode}   & \scriptsize{To ConNode}   & \scriptsize{To   ConNode} & \scriptsize{To  ConNode}  \\
\scriptsize{\textbf{B1}: h-pnt coord. \unitbrackets{m}} & \scriptsize{0.00}         & \scriptsize{0.00}         & \scriptsize{0.00}         & \scriptsize{0.00}         & \scriptsize{0.00}         & \scriptsize{0.00}         \\ 
                                                        & \scriptsize{2500.00}      & \scriptsize{2500.00}      & \scriptsize{2500.00}      & \scriptsize{2500.00}      & \scriptsize{2500.00}      & \scriptsize{2500.00}      \\ 
                                                        & \scriptsize{5000.00}      & \scriptsize{5000.00}      & \scriptsize{5000.00}      & \scriptsize{5000.00}      & \scriptsize{5000.00}      & \scriptsize{5000.00}      \\ 
                                                        & \scriptsize{7500.00}      & \scriptsize{7500.00}      & \scriptsize{7500.00}      & \scriptsize{7500.00}      & \scriptsize{7500.00}      & \scriptsize{7500.00}      \\ 
                                                        & \scriptsize{10000.00}     & \scriptsize{10000.00}     & \scriptsize{10000.00}     & \scriptsize{10000.00}    & \scriptsize{10000.00}     & \scriptsize{10000.00}     \\ 
                                                        & \scriptsize{12500.00}     & \scriptsize{12500.00}     & \scriptsize{12500.00}     & \scriptsize{12500.00}    & \scriptsize{12500.00}     & \scriptsize{12500.00}     \\ 
                                                        & \scriptsize{15000.00}     & \scriptsize{15000.00}     & \scriptsize{15000.00}     & \scriptsize{15000.00}    & \scriptsize{15000.00}     & \scriptsize{15000.00}     \\ 
                                                        & \scriptsize{17500.00}     & \scriptsize{17500.00}     & 											    & 
			 & 											     &                           \\ 
                                                        & \scriptsize{20000.00}     & \scriptsize{20000.00}     & 										    	& 							    
			 & 											     &                            \\ 
                                                        & \scriptsize{22500.00}     & \scriptsize{22500.00}     &                           &              
			 &                           &                           \\ 
                                                        & \scriptsize{25000.00}     & \scriptsize{25000.00}     &                           &               
			 &                           &                           \\ 																			
\hline
\scriptsize{\textbf{B2}: length \unitbrackets{m}}       & \scriptsize{n.a.}         & \scriptsize{n.a.}         & \scriptsize{10000.00}     & \scriptsize{10000.00}     & \scriptsize{10000.00}     & \scriptsize{10000.00}     \\
\scriptsize{\textbf{B2}: direction }                    & \scriptsize{n.a.}         & \scriptsize{n.a.}         & \scriptsize{To ConNode}   & \scriptsize{To ConNode}   & \scriptsize{From ConNode} & \scriptsize{From ConNode} \\
\scriptsize{\textbf{B2}: h-pnt coord. \unitbrackets{m}} & \scriptsize{n.a.}         & \scriptsize{n.a.}         & \scriptsize{0.00}         & \scriptsize{0.00}       & \scriptsize{0.00}         & \scriptsize{0.00}         \\ 
                                                        &                           &                           & \scriptsize{2500.00}      & \scriptsize{2500.00}      & \scriptsize{2500.00}      & \scriptsize{2500.00}      \\
                                                        &                           &                           & \scriptsize{5000.00}      & \scriptsize{5000.00}      & \scriptsize{5000.00}      & \scriptsize{5000.00}      \\ 
                                                        &                           &                           & \scriptsize{7500.00}      & \scriptsize{7500.00}      & \scriptsize{7500.00}      & \scriptsize{7500.00}      \\                                                       
																												&                           &                           & \scriptsize{10000.00}     & \scriptsize{10000.00}    & \scriptsize{10000.00}     & \scriptsize{10000.00}      \\  																										
\hline
\scriptsize{\textbf{B3}: length \unitbrackets{m}}       & \scriptsize{n.a.}         & \scriptsize{n.a.}         & \scriptsize{10000.00}     & \scriptsize{10000.00}     & \scriptsize{10000.00}     & \scriptsize{10000.00}     \\
\scriptsize{\textbf{B3}: direction }                    & \scriptsize{n.a.}         & \scriptsize{n.a.}         & \scriptsize{To ConNode}   & \scriptsize{To ConNode}   & \scriptsize{From ConNode} & \scriptsize{From ConNode} \\
\scriptsize{\textbf{B3}: h-pnt coord. \unitbrackets{m}} & \scriptsize{n.a.}         & \scriptsize{n.a.}         & \scriptsize{0.00}         & \scriptsize{0.00}       & \scriptsize{0.00}         & \scriptsize{0.00}         \\ 
                                                        &                           &                           & \scriptsize{2500.00}      & \scriptsize{2500.00}      & \scriptsize{2500.00}      & \scriptsize{2500.00}      \\
                                                        &                           &                           & \scriptsize{5000.00}      & \scriptsize{5000.00}      & \scriptsize{5000.00}      & \scriptsize{5000.00}      \\ 
                                                        &                           &                           & \scriptsize{7500.00}      & \scriptsize{7500.00}      & \scriptsize{7500.00}      & \scriptsize{7500.00}      \\                                                       
																												&                           &                           & \scriptsize{10000.00}     & \scriptsize{10000.00}    & \scriptsize{10000.00}     & \scriptsize{10000.00}      \\  																					
[1ex] \hline
\end{longtable}



\begin{longtable}{|p{13.5mm}|p{10.5mm}|p{16.0mm}|p{29.5mm}|p{47.5mm}|}
\caption{Boundary conditions applied in sub-models \textcolor[rgb]{1,0,0}{\textbf{T1}} to \textcolor[rgb]{1,0,0}{\textbf{T6}}.} 
\label{10407:Boundary conditions} \\ 
\hline
     \begin{center} \scriptsize\textbf{Sub-model} \end{center}
   & \begin{center} \scriptsize\textbf{Branch} \end{center} 
	 & \begin{center} \scriptsize\textbf{x-coord. \unitbrackets{m}} \end{center} 
	 & \begin{center} \scriptsize\textbf{ID} \end{center} 
	& \begin{center} \scriptsize\textbf{Boundary condition} \end{center} \\ 
[1ex] \hline \endhead
\hline
\endfoot
%
\STRUT
\scriptsize\textcolor[rgb]{1,0,0}{\textbf{T1}} & \scriptsize\textbf{B1} & \scriptsize{0.00}     & \scriptsize{T1\_h-Bnd\_B1\_x0m}     & \scriptsize{Constant water level of 2.50 m AD}                                                \\
                                               & \scriptsize\textbf{B1} & \scriptsize{25000.00} & \scriptsize{T1\_Q-Bnd\_B1\_x25000m} & \scriptsize{500 times discharge, depicted in \autoref{fig:10407:Upstream discharge hydrograph.jpg}} \\
\hline
\scriptsize\textcolor[rgb]{1,0,0}{\textbf{T2}} & \scriptsize\textbf{B1} & \scriptsize{25000.00} & \scriptsize{T2\_h-Bnd\_B1\_x25000m} & \scriptsize{Constant water level of 2.50 m AD}                                                \\
                                               & \scriptsize\textbf{B1} & \scriptsize{0.00}     & \scriptsize{T2\_Q-Bnd\_B1\_x0m}     & \scriptsize{500 times discharge, depicted in \autoref{fig:10407:Upstream discharge hydrograph.jpg}} \\
\hline
\scriptsize\textcolor[rgb]{1,0,0}{\textbf{T3}} & \scriptsize\textbf{B1} & \scriptsize{0.00}     & \scriptsize{T3\_h-Bnd\_B1\_x0m}     & \scriptsize{Constant water level of 2.50 m AD}                                                \\
                                               & \scriptsize\textbf{B2} & \scriptsize{0.00}     & \scriptsize{T3\_Q-Bnd\_B2\_x0m}     & \scriptsize{50 times discharge, depicted in \autoref{fig:10407:Upstream discharge hydrograph.jpg}}  \\
                                               & \scriptsize\textbf{B3} & \scriptsize{0.00}     & \scriptsize{T3\_Q-Bnd\_B3\_x0m}     & \scriptsize{450 times discharge, depicted in \autoref{fig:10407:Upstream discharge hydrograph.jpg}} \\
\hline
\scriptsize\textcolor[rgb]{1,0,0}{\textbf{T4}} & \scriptsize\textbf{B1} & \scriptsize{0.00}     & \scriptsize{T4\_h-Bnd\_B1\_x0m}     & \scriptsize{Constant water level of 2.50 m AD}                                                \\
                                               & \scriptsize\textbf{B2} & \scriptsize{0.00}     & \scriptsize{T4\_Q-Bnd\_B2\_x0m}     & \scriptsize{100 times discharge, depicted in \autoref{fig:10407:Upstream discharge hydrograph.jpg}} \\
                                               & \scriptsize\textbf{B3} & \scriptsize{10000.00} & \scriptsize{T4\_Q-Bnd\_B3\_x10000m} & \scriptsize{400 times discharge, depicted in \autoref{fig:10407:Upstream discharge hydrograph.jpg}} \\
\hline
\scriptsize\textcolor[rgb]{1,0,0}{\textbf{T5}} & \scriptsize\textbf{B1} & \scriptsize{0.00}     & \scriptsize{T5\_h-Bnd\_B1\_x0m}     & \scriptsize{Constant water level of 2.50 m AD}                                                \\
                                               & \scriptsize\textbf{B2} & \scriptsize{10000.00} & \scriptsize{T5\_Q-Bnd\_B2\_x10000m} & \scriptsize{150 times discharge, depicted in \autoref{fig:10407:Upstream discharge hydrograph.jpg}} \\
                                               & \scriptsize\textbf{B3} & \scriptsize{0.00}     & \scriptsize{T5\_Q-Bnd\_B3\_x0m}     & \scriptsize{350 times discharge, depicted in \autoref{fig:10407:Upstream discharge hydrograph.jpg}} \\
\hline
\scriptsize\textcolor[rgb]{1,0,0}{\textbf{T6}} & \scriptsize\textbf{B1} & \scriptsize{0.00}     & \scriptsize{T6\_h-Bnd\_B1\_x0m}     & \scriptsize{Constant water level of 2.50 m AD}                                                \\
                                               & \scriptsize\textbf{B2} & \scriptsize{10000.00} & \scriptsize{T6\_Q-Bnd\_B2\_x10000m} & \scriptsize{250 times discharge, depicted in \autoref{fig:10407:Upstream discharge hydrograph.jpg}} \\
                                               & \scriptsize\textbf{B3} & \scriptsize{10000.00} & \scriptsize{T6\_Q-Bnd\_B3\_x10000m} & \scriptsize{250 times discharge, depicted in \autoref{fig:10407:Upstream discharge hydrograph.jpg}} \\
[1ex] \hline
\end{longtable}


\begin{longtable}{|p{13.5mm}|p{10.5mm}|p{16.0mm}|p{26.0mm}|p{12.5mm}|p{16.0mm}|p{19.5mm}|}
\caption{"Open Y-Z rectangular cross-sections (using the vertically segmented conveyance method)" in sub-model \textcolor[rgb]{1,0,0}{\textbf{T1}} to \textcolor[rgb]{1,0,0}{\textbf{T6}}.} 
\label{10407:Open Y-Z rectangular cross-sections} \\ 
\hline
     \begin{center} \scriptsize\textbf{Sub-model} \end{center}
   & \begin{center} \scriptsize\textbf{Branch} \end{center} 
	 & \begin{center} \scriptsize\textbf{x-coord. \unitbrackets{m}} \end{center} 
	 & \begin{center} \scriptsize\textbf{ID} \end{center} 
	 & \begin{center} \scriptsize\textbf{Width \unitbrackets{m}} \end{center}
	 & \begin{center} \scriptsize\textbf{Bed level \unitbrackets{m}} \end{center}
	& \begin{center} \scriptsize\textbf{Chezy \unitbrackets{$m^{0.5}/s$}} \end{center} \\ 
[1ex] \hline \endhead
\hline
\endfoot
%
\STRUT
\scriptsize\textcolor[rgb]{1,0,0}{\textbf{T1}} & \scriptsize\textbf{B1} & \scriptsize{0.00}     & \scriptsize{T1\_Crs\_B1\_x0m}     & \scriptsize{500.00}  & \scriptsize{0.00}  & \scriptsize{50.0} \\
                                               & \scriptsize\textbf{B1} & \scriptsize{15000.00} & \scriptsize{T1\_Crs\_B1\_x15000m} & \scriptsize{500.00}  & \scriptsize{10.00} & \scriptsize{50.0} \\
                                               & \scriptsize\textbf{B1} & \scriptsize{25000.00} & \scriptsize{T1\_Crs\_B1\_x25000m} & \scriptsize{500.00}  & \scriptsize{10.15} & \scriptsize{50.0} \\
\hline
\scriptsize\textcolor[rgb]{1,0,0}{\textbf{T2}} & \scriptsize\textbf{B1} & \scriptsize{25000.00} & \scriptsize{T2\_Crs\_B1\_x25000m} & \scriptsize{500.00}  & \scriptsize{0.00}  & \scriptsize{50.0} \\
                                               & \scriptsize\textbf{B1} & \scriptsize{10000.00} & \scriptsize{T2\_Crs\_B1\_x10000m} & \scriptsize{500.00}  & \scriptsize{10.00} & \scriptsize{50.0} \\
                                               & \scriptsize\textbf{B1} & \scriptsize{0.00}     & \scriptsize{T2\_Crs\_B1\_x0m}     & \scriptsize{500.00}  & \scriptsize{10.15} & \scriptsize{50.0} \\
\hline
\scriptsize\textcolor[rgb]{1,0,0}{\textbf{T3}} & \scriptsize\textbf{B1} & \scriptsize{0.00}     & \scriptsize{T3\_Crs\_B1\_x0m}     & \scriptsize{500.00}  & \scriptsize{0.00}  & \scriptsize{50.0} \\
                                               & \scriptsize\textbf{B1} & \scriptsize{15000.00} & \scriptsize{T3\_Crs\_B1\_x15000m} & \scriptsize{500.00}  & \scriptsize{10.00} & \scriptsize{50.0} \\
                                               & \scriptsize\textbf{B2} & \scriptsize{10000.00} & \scriptsize{T3\_Crs\_B2\_x10000m} & \scriptsize{50.00}   & \scriptsize{10.00} & \scriptsize{50.0} \\
                                               & \scriptsize\textbf{B2} & \scriptsize{0.00}     & \scriptsize{T3\_Crs\_B2\_x0m}     & \scriptsize{50.00}   & \scriptsize{10.15} & \scriptsize{50.0} \\
                                               & \scriptsize\textbf{B3} & \scriptsize{10000.00} & \scriptsize{T3\_Crs\_B3\_x10000m} & \scriptsize{450.00}  & \scriptsize{10.00} & \scriptsize{50.0} \\
                                               & \scriptsize\textbf{B3} & \scriptsize{0.00}     & \scriptsize{T3\_Crs\_B3\_x0m}     & \scriptsize{450.00}  & \scriptsize{10.15} & \scriptsize{50.0} \\
\hline
\scriptsize\textcolor[rgb]{1,0,0}{\textbf{T4}} & \scriptsize\textbf{B1} & \scriptsize{0.00}     & \scriptsize{T4\_Crs\_B1\_x0m}     & \scriptsize{500.00}  & \scriptsize{0.00}  & \scriptsize{50.0} \\
                                               & \scriptsize\textbf{B1} & \scriptsize{15000.00} & \scriptsize{T4\_Crs\_B1\_x15000m} & \scriptsize{500.00}  & \scriptsize{10.00} & \scriptsize{50.0} \\
                                               & \scriptsize\textbf{B2} & \scriptsize{10000.00} & \scriptsize{T4\_Crs\_B2\_x10000m} & \scriptsize{100.00}  & \scriptsize{10.00} & \scriptsize{50.0} \\
                                               & \scriptsize\textbf{B2} & \scriptsize{0.00}     & \scriptsize{T4\_Crs\_B2\_x0m}     & \scriptsize{100.00}  & \scriptsize{10.15} & \scriptsize{50.0} \\
                                               & \scriptsize\textbf{B3} & \scriptsize{0.00}     & \scriptsize{T4\_Crs\_B3\_x0m}     & \scriptsize{400.00}  & \scriptsize{10.00} & \scriptsize{50.0} \\
                                               & \scriptsize\textbf{B3} & \scriptsize{10000.00} & \scriptsize{T4\_Crs\_B3\_x10000m} & \scriptsize{400.00}  & \scriptsize{10.15} & \scriptsize{50.0} \\
\hline
\scriptsize\textcolor[rgb]{1,0,0}{\textbf{T5}} & \scriptsize\textbf{B1} & \scriptsize{0.00}     & \scriptsize{T5\_Crs\_B1\_x0m}     & \scriptsize{500.00}  & \scriptsize{0.00}  & \scriptsize{50.0} \\
                                               & \scriptsize\textbf{B1} & \scriptsize{15000.00} & \scriptsize{T5\_Crs\_B1\_x15000m} & \scriptsize{500.00}  & \scriptsize{10.00} & \scriptsize{50.0} \\
                                               & \scriptsize\textbf{B2} & \scriptsize{0.00}     & \scriptsize{T5\_Crs\_B2\_x0m}     & \scriptsize{150.00}  & \scriptsize{10.00} & \scriptsize{50.0} \\
                                               & \scriptsize\textbf{B2} & \scriptsize{10000.00} & \scriptsize{T5\_Crs\_B2\_x10000m} & \scriptsize{150.00}  & \scriptsize{10.15} & \scriptsize{50.0} \\
                                               & \scriptsize\textbf{B3} & \scriptsize{10000.00} & \scriptsize{T5\_Crs\_B3\_x10000m} & \scriptsize{350.00}  & \scriptsize{10.00} & \scriptsize{50.0} \\
                                               & \scriptsize\textbf{B3} & \scriptsize{0.00}     & \scriptsize{T5\_Crs\_B3\_x0m}     & \scriptsize{350.00}  & \scriptsize{10.15} & \scriptsize{50.0} \\
\hline
\scriptsize\textcolor[rgb]{1,0,0}{\textbf{T6}} & \scriptsize\textbf{B1} & \scriptsize{0.00}     & \scriptsize{T6\_Crs\_B1\_x0m}     & \scriptsize{500.00}  & \scriptsize{0.00}  & \scriptsize{50.0} \\
                                               & \scriptsize\textbf{B1} & \scriptsize{15000.00} & \scriptsize{T6\_Crs\_B1\_x15000m} & \scriptsize{500.00}  & \scriptsize{10.00} & \scriptsize{50.0} \\
                                               & \scriptsize\textbf{B2} & \scriptsize{0.00}     & \scriptsize{T6\_Crs\_B2\_x0m}     & \scriptsize{250.00}  & \scriptsize{10.00} & \scriptsize{50.0} \\
                                               & \scriptsize\textbf{B2} & \scriptsize{10000.00} & \scriptsize{T6\_Crs\_B2\_x10000m} & \scriptsize{250.00}  & \scriptsize{10.15} & \scriptsize{50.0} \\
                                               & \scriptsize\textbf{B3} & \scriptsize{0.00}     & \scriptsize{T6\_Crs\_B3\_x0m}     & \scriptsize{250.00}  & \scriptsize{10.00} & \scriptsize{50.0} \\
                                               & \scriptsize\textbf{B3} & \scriptsize{10000.00} & \scriptsize{T6\_Crs\_B3\_x10000m} & \scriptsize{250.00}  & \scriptsize{10.15} & \scriptsize{50.0} \\
[1ex] \hline
\end{longtable}


%---------------------------------------------------------------------------------
\section*{Results}

Water level and velocity results are shown for sub-models \textbf{\textcolor[rgb]{1,0,0}{T1}} to \textbf{\textcolor[rgb]{1,0,0}{T3}} on \autoref{fig:10407:FMWaterlevelsAtNodes} to \autoref{fig:10407:FMWaterVelT3B2}.

\begin{figure}[H]
    \centering
    \includegraphics*[width=0.85\textwidth]{figures/e02_f104_c07_FM_Water_Level.jpg}
    \caption{Water levels at T3\_ConNode and T1\_h\_pnt\_x15000m (see \autoref{fig:10407:Test model lay-out}) for applying junction advection. Water level differences (T1\_h\_pnt\_x15000m \textcolor[rgb]{1,0,0}{minus} T3\_ConNode) are plotted against the right vertical axis.}
    \label{fig:10407:FMWaterlevelsAtNodes}
\end{figure} 

\begin{figure}[H]
    \centering
    \includegraphics*[width=0.85\textwidth]{figures/e02_f104_c07_FM_Velocity_T1_T3B1.jpg}
    \caption{Velocities at "u-pnt on B1 West of T1\_h-pnt\_x15000m" and "u-pnt on B1 next to T3\_ConNode" (see \autoref{fig:10407:Test model lay-out}) for applying junction advection. Differences in velocities ("u-pnt on B1 West of T1\_h-pnt\_x15000m" \textcolor[rgb]{1,0,0}{minus} "u-pnt on B1 next to T3\_ConNode") are plotted against the right vertical axis.}
    \label{fig:10407:FMWaterVelT3B1}
\end{figure} 

\begin{figure}[H]
    \centering
    \includegraphics*[width=0.85\textwidth]{figures/e02_f104_c07_FM_Velocity_T2_T3B2.jpg}
    \caption{Velocities at "u-pnt on B1 East of T2\_h-pnt\_x10000m" and "u-pnt on B2 next to T3\_ConNode" (see \autoref{fig:10407:Test model lay-out}) for applying junction advection. Differences in velocities ("u-pnt on B1 East of T2\_h-pnt\_x10000m" \textcolor[rgb]{1,0,0}{minus} "u-pnt on B2 next to T3\_ConNode")  are plotted against the right vertical axis.}
    \label{fig:10407:FMWaterVelT3B2}
\end{figure} 

Water level and velocity results for sub-models \textbf{\textcolor[rgb]{1,0,0}{T1}} to \textbf{\textcolor[rgb]{1,0,0}{T6}} ran on DFlowFM  are compared to the results of a similar test case ran on SOBEK3 . These are shown on \autoref{fig:10407:SBKFMWaterLevel} to \autoref{fig:10407:SBKFMWaterVelTable}.

\begin{figure}[H]
    \centering
    \includegraphics*[width=0.85\textwidth]{figures/e02_f104_c07_SOBEK_FM_Water_Level.jpg}
    \caption{Water level results at T1\_h\_pnt\_x15000m (see \autoref{fig:10407:Test model lay-out}) calculated through DFlowFM and SOBEK3. Water level differences (DFlowFM \textcolor[rgb]{1,0,0}{minus} SOBEK3) are plotted against the right vertical axis.}
    \label{fig:10407:SBKFMWaterLevel}
\end{figure} 

\begin{figure}[H]
    \centering
    \includegraphics*[width=0.85\textwidth]{figures/e02_f104_c07_SOBEK_FM_Velocity_B1.jpg}
    \caption{Velocities at "u-pnt on B1 next to T3\_ConNode" (see \autoref{fig:10407:Test model lay-out}) calculated through DFlowFM and SOBEK3. Differences in velocities (DFlowFM \textcolor[rgb]{1,0,0}{minus} SOBEK3) are plotted against the right vertical axis.}
    \label{fig:10407:SBKFMWaterVelB1}
\end{figure} 

\begin{figure}[H]
    \centering
    \includegraphics*[width=0.85\textwidth]{figures/e02_f104_c07_SOBEK_FM_Velocity_B2.jpg}
    \caption{Velocities at "u-pnt on B2 next to T3\_ConNode" (see \autoref{fig:10407:Test model lay-out}) calculated through DFlowFM and SOBEK3. Differences in velocities (DFlowFM \textcolor[rgb]{1,0,0}{minus} SOBEK3) are plotted against the right vertical axis.}
    \label{fig:10407:SBKFMWaterVelB2}
\end{figure} 

\begin{figure}[H]
    \centering
    \includegraphics*[width=0.85\textwidth]{figures/e02_f104_c07_SOBEK_FM_Water_Level_table.jpg}
    \caption{Minimum and maximum differences (DFlowFM \textcolor[rgb]{1,0,0}{minus} SOBEK3) in water levels at (connection) nodes, corresponding to the time period from 02-01-2015 00:00:00 to 04-01-2015 00:00:00.}
    \label{fig:10407:SBKFMWaterLevelTable}
\end{figure} 

\begin{figure}[H]
    \centering
    \includegraphics*[width=0.85\textwidth]{figures/e02_f104_c07_SOBEK_FM_Velocity_table.jpg}
    \caption{Minimum and maximum differences in velocities (DFlowFM \textcolor[rgb]{1,0,0}{minus} SOBEK3), corresponding to the time period from 02-01-2015 00:00:00 to 04-01-2015 00:00:00.}
    \label{fig:10407:SBKFMWaterVelTable}
\end{figure} 


%---------------------------------------------------------------------------------
\section*{Analysis of results}

The results are analyzed through (1) a comparison of the results for sub-models with (i.e. \textbf{\textcolor[rgb]{1,0,0}{T3}}) and without (i.e. \textbf{\textcolor[rgb]{1,0,0}{T1}} and \textbf{\textcolor[rgb]{1,0,0}{T2}}) junctions (or connection nodes) and (2) a comparison of the results to an identical SOBEK3 case. 

As described in \autoref{10407:model lay-out}, per metre of cross-sectional width yields that the hydraulic forcing, the bathymetry and roughness of sub-models \textbf{\textcolor[rgb]{1,0,0}{T3}}, \textbf{\textcolor[rgb]{1,0,0}{T4}}, \textbf{\textcolor[rgb]{1,0,0}{T5}} and \textbf{\textcolor[rgb]{1,0,0}{T6}} are identical to those of sub-models \textbf{\textcolor[rgb]{1,0,0}{T1}} and \textbf{\textcolor[rgb]{1,0,0}{T2}}. Hence, if junction advection (see \autoref{10407:description of junction advection}) is applied, computed water levels and flow velocities in sub-models \textbf{\textcolor[rgb]{1,0,0}{T3}} to \textbf{\textcolor[rgb]{1,0,0}{T6}} should be identical to those computed in sub-models \textbf{\textcolor[rgb]{1,0,0}{T1}} and \textbf{\textcolor[rgb]{1,0,0}{T2}}.

From a comparison of the results for sub-models with (i.e. \textbf{\textcolor[rgb]{1,0,0}{T3}}) and without (i.e. \textbf{\textcolor[rgb]{1,0,0}{T1}} and \textbf{\textcolor[rgb]{1,0,0}{T2}}) junctions (or connection nodes) on \autoref{fig:10407:FMWaterlevelsAtNodes} to \autoref{fig:10407:FMWaterVelT3B2} it is observed that the water level and velocity calculated in sub-model  \textbf{\textcolor[rgb]{1,0,0}{T3}} is almost identical to those calculated in sub-models \textbf{\textcolor[rgb]{1,0,0}{T1}} and \textbf{\textcolor[rgb]{1,0,0}{T2}}. The evident differences are insignificant and do not cause further numerical disturbances.

From a comparison of the results to an identical SOBEK3 case it is observed that the water levels and velocities for all sub-models are almost identical. The differences observed may be partly attributed to:
\begin{enumerate}
 \item the different initial conditions that are applied in the SOBEK3 case. This statement is supported by the differences of similar magnitude evident also in the sub-models without junction advection applied,  \textbf{\textcolor[rgb]{1,0,0}{T1}} and \textbf{\textcolor[rgb]{1,0,0}{T2}};
 \item the directional change at the branches in line with the findings of test case \textbf{f115}\_numerical\_aspects\_\textbf{c16}\_curved-channel-junction-advection-reference); and
 \item a difference in the distribution of flow after a bifurcation which causes a difference in the velocities. The magnitude of the difference seems to be inversely related to the width of the branch. The largest difference is observed on Branch B2 of sub-model  \textbf{\textcolor[rgb]{1,0,0}{T3}} with the smallest width (50m) followed by B2 of sub-model  \textbf{\textcolor[rgb]{1,0,0}{T4}} with a width of 100m.
\end{enumerate}

It is not understood why there is a general tendency for DFlowFM to predict lower water levels [and lower velocities] at the rise of the flood wave and higher water levels [and higher velocities] at the fall of the flood wave.

These observed differences however are insignificant and do not cause further numerical disturbances. 


%---------------------------------------------------------------------------------
\subsection{Junction advection} \label{10407:description of junction advection}
Junction advection (or accounting for the advection term ($u \pdiff{u}{x}$) near connection nodes) is incorporated in the dflowFM1d (cf\_dll.dll) computational core in accordance with \citet{Borsboom2017}.

Junction advection is not applied at boundary nodes. For connection nodes yields that junction advection is only applied if more than two branches are connected at a connection node. The reason, herefore, are:
\begin{itemize}
	\item For a branch connected at a boundary node, the advection term ($u \pdiff{u}{x}$) near the boundary node is determined by taking the velocity ($u$) equal to zero at the boundary node.
	\item For connection node having only one branch connected to it, the advection term ($u \pdiff{u}{x}$) near the connection node is determined by taking the velocity ($u$) equal to zero at the connection node.	
	\item For two connected branches, the connection node is just treated as an ordinary water-level point. In other words the advection term ($u \pdiff{u}{x}$) is computed in the same manner as for any water-level point, located in a branch at sufficient distance of a connection node or boundary node.
\end{itemize} 

For connection nodes having \underline{more than two} connected branches, junction advection means that the advection term ($u \pdiff{u}{x}$) is taken into account in solving the St. Venant momentum equation between a connection node (or junction) and the first adjacent water-level-point, located on a particular connected branch.

Before junction advection was implemented (i.e. no junction advection was applied) in the 'dflow1d' computational core following was done. For connection nodes having more than two connected branches, the advection term ($u \pdiff{u}{x}$) directly adjacent to such connection node was set equal to zero in solving the St. Venant momentum equation. 

\Note Junction advection does \textbf{\textcolor[rgb]{1,0,0}{not}} account for (possible) hydraulic head-losses (or \textbf{\textcolor[rgb]{1,0,0}{junction-losses}}) at a connection node. The reason being that junction advection does \underline{not} take into account the angles (in the horizontal and/or vertical plane) between the various branches, that are connected at a connection node (or junction). 

%---------------------------------------------------------------------------------
\printrefsegment
